\subsection{Multi-Modal Learning}
% Multi-modal learning aims to build models that can process and relate information from multiple modalities~\citep{baltruvsaitis2018multimodal}. 
% Vision-language tasks, such as VQA~\citep{antol2015vqa, gao2019dynamic}, image captioning~\citep{xu2015show}, require an understanding of both textual and visual semantics. 
Most vision-language models~\cite{chen2020uniter,lu2019vilbert,li2020oscar} use a bunch of cross-modal transformers to fuse and align the information between text and image.
% Most vision-language models use a bunch of cross-modal transformers to fuse and align the information between text and image, such as LXMERT~\citep{tan2019lxmert}, UNITER~\citep{chen2020uniter}, ViLBERT~\citep{lu2019vilbert}, VisualBERT~\citep{li2019visualbert}, OSCAR~\citep{li2020oscar}, Pixel-BERT~\citep{huang2020pixel}. 
These methods either need an off-the-shelf object detector to extract region features or dedicated cross-modal transformer layers, significantly hindering their scalability. 
Our \workname, by contrast, uses a simple yet effective two-tower framework with multi-modal interaction only at the top. Moreover, this series of models~\citep{radford2021learning, jia2021scaling, huo2021wenlan} can perform zero-shot recognition, adapting to new categories with no seen labeled data. ~\citet{shen2021much} also shows that the pre-trained CLIP model can significantly benefit the downstream VQA and image caption tasks.
Our \workname is supposed to be compatible with more modalities, e.g., acoustic signals~\citep{akbari2021vatt}. More modalities included, more correlated supervision are expected to be exploited. 
% We hope our \workname can inspire and advance the future study of multi-modal learning.
